\documentclass[../../../include/open-logic-section]{subfiles}

\begin{document}

\olfileid{sth}{story}{rus}
\olsection{रसेलची विरोधापत्ती (पुन्हा)}

\olref[sfr][][]{part}
मध्ये आपण अनौपचारिक
संचसिद्धान्त वापरला.
पण एका \emph{अतिशय}
भोळ्या कल्पनेनुसार
संच म्हणजे विधेयांचे
विस्तारच असतात.
या कल्पनेतून पुढील
तत्त्व स्वीकारावे लागेल:

\begin{defish}
  \emph{अनौपचारिक
  गुणधर्माधारित संचरचना.}
  कोणत्याही सूत्र~$\phi$
  साठी
  $\Setabs{x}{\phi(x)}$
  अस्तित्वात असतो.
\end{defish}

हे तत्त्व मोहक
असले, तरी ते
सिद्धपणे विसंगत
आहे. \olref[sfr][set][rus]{sec}
मध्ये आपण हे
पाहिले आहे.
पण निष्कर्ष
इतका महत्त्वाचा
आणि इतका
सरळ आहे
की तो
शब्दशः पुन्हा
मांडण्यासारखा आहे.

\begin{thm}[रसेलची विरोधापत्ती]
$R = \Setabs{x}{x \notin x}$
असा कोणताही संच नाही.
\end{thm}

\begin{proof}
$R = \Setabs{x}{x \notin x}$
अस्तित्वात असेल, तर
$R \in R$ तेव्हा
आणि केवळ तेव्हाच
$R \notin R$;
हा व्याघात आहे.
\end{proof}

रसेल यांनी हा
निष्कर्ष जून १९०१
मध्ये शोधला. (मात्र
त्यांनी विरोधापत्ती
आपण आत्ताच
मांडलेल्या रूपात
दिलेली नव्हती,
कारण ते
\emph{Grundgesetze}
मध्ये मांडलेल्या
फ्रेग यांच्या
संच उपपत्तीचा
विचार करत होते.
\olref[blv]{sec}
मध्ये आपण
याकडे परत
येऊ.) फ्रेग
यांच्या प्रणालीतील
विसंगती स्पष्ट
करण्यासाठी रसेल
यांनी त्यांना
१६ जून १९०२
रोजी पत्र
लिहिले. त्या
पत्रव्यवहारासाठी
आणि थोड्या
पार्श्वभूमीसाठी
\citet[pp.~124--8]{Heijenoort1967}
पाहा.

या दोन ओळींची
सिद्धता \emph{शुद्ध
तर्कशास्त्रातून}
मिळते, यावर
भर द्यायला हवा.
मान्य आहे की
$\Setabs{x}{x \notin x}$
या लेखनात
आपण विस्तारात्मकतेचे
एक (तार्किकेतर?)
स्वयंसिद्धक अप्रत्यक्षपणे
वापरले आहे:
$\Setabs{x}{\phi(x)}$
म्हणजे, असा संच
असेल तर,
$\phi$ पूर्ण
करणाऱ्या वस्तूंचा
\emph{एकमेव}
(विस्तारात्मकतेमुळे)
संच. पण निष्कर्ष
पुढीलप्रमाणे
मांडल्यास
विस्तारात्मकतेचा
आभासही टाळता
येतो: \emph{ज्या
संचांचे स्वतःत
सदस्यत्व नाही,
नेमके तेच
ज्याचे सदस्य
आहेत असा
संच नाही}.
याचा संचांशीही
फारसा विशेष
संबंध नाही.
रसेल यांनी
स्वतः नमूद
केल्याप्रमाणे,
अगदी तसाच
युक्तिवाद
दाखवतो:
\emph{जे पुरुष
स्वतःची दाढी
करत नाहीत,
नेमक्या त्याच
पुरुषांची दाढी
करणारा पुरुष
नाही}. किंवा:
\emph{जी पग
जातीची कुत्री
स्वतःचा वास
घेत नाहीत,
नेमक्या त्याच
पग कुत्र्यांचा
वास घेणारा
पग कुत्रा
नाही}. इत्यादी.
सांकेतिक रूपात
हा निष्कर्ष
इतकाच:
\[
\lnot \exists x \forall z(Rzx \liff \lnot R zz).
\]
हे प्रथम-क्रम
तर्कशास्त्राचे
प्रमेय(-रूपबंध)
आहे. त्यामुळे
संच उपपत्तीत
फक्त फेरबदल
करून रसेलची
विरोधापत्ती
टाळता येत
नाही: संच
उपपत्तीपर्यंत
\emph{पोहोचण्याआधीच}
ती उद्भवते.
आपण (अभिजात)
प्रथम-क्रम
तर्कशास्त्र
वापरणार असू,
तर
$R = \Setabs{x}{x\notin x}$
असा संच नाही
हे \emph{स्वीकारावेच}
लागते.

सार असा:
विसंगती
\emph{टाळूनही}
अनौपचारिक
गुणधर्माधारित
संचरचना
स्वीकारायची
असेल, तर
केवळ \emph{संच
उपपत्तीत}
किरकोळ फेरबदल
पुरेसे नाहीत.
त्याऐवजी
\emph{तर्कशास्त्र}
मुळापासून
बदलावे लागेल.

अर्थात,
अभिजातेतर
तर्कशास्त्रे
वापरणाऱ्या
संच उपपत्ती
मांडल्या गेल्या
आहेत. पण
त्या, सौम्य
शब्दांत सांगायचे
तर, मानक
नाहीत.
रसेलच्या
विरोधापत्तीला
सामोरे जाण्याची
मानक पद्धत
म्हणजे तिला
संच अस्तित्वात
नसल्याची
सरळ सिद्धता
मानणे, आणि
मग या निष्कर्षासह
काम कसे करावे
हे शिकणे.
आपण हीच
पद्धत वापरू.

\end{document}
